\documentclass[12pt,a4paper]{article}
\usepackage{expediente}
\begin{document}
% =============================================================================
% FICHA DE DATOS DEL EXPEDIENTE --- actualizado 20 ago 2026
% =============================================================================

% --- VENDEDOR / TITULAR REGISTRAL -------------------------------------------
\newcommand{\VendedorNombres}{RICHAR DARWIN}
\newcommand{\VendedorApellidos}{CUTIRE CONDORI}
\newcommand{\VendedorNombreCompleto}{RICHAR DARWIN CUTIRE CONDORI}
\newcommand{\VendedorDNI}{42610690}
\newcommand{\VendedorEstadoCivil}{soltero, sin unión de hecho vigente}
\newcommand{\VendedorOcupacion}{\faltante{[ocupación de Richar]}}
\newcommand{\VendedorDomicilio}{\faltante{[domicilio actual de Richar]}}
\newcommand{\VendedorTelefono}{\faltante{[teléfono de Richar]}}
\newcommand{\VendedorCorreo}{\faltante{[correo]}}
\newcommand{\VendedorNacionalidad}{peruano}
\newcommand{\VendedorFechaNac}{\faltante{[fecha de nacimiento]}}
\newcommand{\VendedorDNIEmision}{\faltante{[fecha de emisión del DNI de Richar --- la pide SUNARP]}}

% --- CONVIVIENTE: NO CORRESPONDE -------------------------------------------
\newcommand{\ConvivienteNombre}{NO CORRESPONDE}
\newcommand{\ConvivienteDNI}{NO CORRESPONDE}
\newcommand{\ConvivienteDesde}{NO CORRESPONDE}
\newcommand{\ConvivienteInscrita}{no existe unión de hecho vigente}
\newcommand{\ConvivienteDomicilio}{NO CORRESPONDE}

% --- COMPRADORES (copropiedad en partes iguales) ---------------------------
\newcommand{\CompradorUnoNombre}{NICO ALVARO CUTIRE ARCE}
\newcommand{\CompradorUnoDNI}{48325017}
\newcommand{\CompradorUnoEmision}{01 de enero de 2024}
\newcommand{\CompradorUnoEstadoCivil}{soltero}
\newcommand{\CompradorUnoOcupacion}{ingeniero informático}
\newcommand{\CompradorUnoDomicilio}{Jr.\ Wiracocha, Sicuani, provincia de Canchis, Cusco, Perú}
\newcommand{\CompradorUnoTelefono}{984715108}

\newcommand{\CompradorDosNombre}{LUCY CUTIRE ARCE}
\newcommand{\CompradorDosDNI}{47478852}
\newcommand{\CompradorDosEmision}{09 de enero de 2026}
\newcommand{\CompradorDosEstadoCivil}{soltera}
\newcommand{\CompradorDosOcupacion}{contador}
\newcommand{\CompradorDosDomicilio}{Comunidad Hanocca, distrito de Layo, provincia de Canas, Cusco, Perú}
\newcommand{\CompradorDosTelefono}{953440875}

\newcommand{\CompradorNombres}{NICO ALVARO y LUCY}
\newcommand{\CompradorApellidos}{CUTIRE ARCE}
\newcommand{\CompradorNombreCompleto}{NICO ALVARO CUTIRE ARCE y LUCY CUTIRE ARCE}
\newcommand{\CompradorDNI}{48325017 y 47478852}
\newcommand{\CompradorEstadoCivil}{solteros}
\newcommand{\CompradorConyuge}{NO CORRESPONDE}
\newcommand{\CompradorConyugeDNI}{NO CORRESPONDE}
\newcommand{\CompradorOcupacion}{ingeniero informático y contador, respectivamente}
\newcommand{\CompradorDomicilio}{Sicuani (Canchis) y Layo (Canas), Cusco}
\newcommand{\CompradorTelefono}{984715108 / 953440875}
\newcommand{\CompradorCorreo}{\faltante{[correo de contacto]}}
\newcommand{\CompradorNacionalidad}{peruanos}
\newcommand{\PorcentajeCopropiedad}{50\,\% cada uno}

% --- INTERMEDIARIO / OCUPANTE DEL REMANENTE --------------------------------
\newcommand{\IntermediarioNombre}{LUCIO GIL CUTIRE PARI}
\newcommand{\IntermediarioDNI}{42467850}
\newcommand{\IntermediarioEstadoCivil}{\faltante{[estado civil de Lucio]}}
\newcommand{\IntermediarioOcupacion}{\faltante{[ocupación]}}
\newcommand{\IntermediarioDomicilio}{\faltante{[domicilio de Lucio]}}
\newcommand{\IntermediarioTelefono}{\faltante{[teléfono de Lucio]}}
\newcommand{\IntermediarioParentesco}{familiar (apellido Cutire)}

% --- PREDIO MATRIZ ----------------------------------------------------------
\newcommand{\PartidaMatriz}{\faltante{[N.º de partida electrónica SUNARP --- pendiente]}}
\newcommand{\OficinaRegistral}{Oficina Registral de Puerto Maldonado}
\newcommand{\AreaMatriz}{\faltante{[área total inscrita; las partes refieren del orden de 100 ha]}}
\newcommand{\UnidadCatastral}{\faltante{[código catastral / UC, si existe]}}
\newcommand{\NombreFundo}{\faltante{[nombre del fundo, si tiene]}}
\newcommand{\KmCarretera}{Interoceánica Sur, al este del Puente Jayave (el puente queda al oeste del predio, sobre el río; no debajo del lote)}

% --- PREDIO A INDEPENDIZAR --------------------------------------------------
\newcommand{\AreaIndependizarHa}{10{,}00}
\newcommand{\AreaIndependizarMetros}{100 000{,}00}
\newcommand{\PerimetroM}{2 200{,}58}
\newcommand{\FrenteM}{100{,}29}
\newcommand{\FondoM}{1 000{,}00}

\newcommand{\PUnoLat}{$-12{,}912452$}
\newcommand{\PUnoLon}{$-70{,}170058$}
\newcommand{\PDosLat}{$-12{,}912497$}
\newcommand{\PDosLon}{$-70{,}170981$}
\newcommand{\PTresLat}{$-12{,}903469$}
\newcommand{\PTresLon}{$-70{,}171438$}
\newcommand{\PCuatroLat}{$-12{,}903424$}
\newcommand{\PCuatroLon}{$-70{,}170515$}

\newcommand{\PUnoE}{373 065{,}00}
\newcommand{\PUnoN}{8 572 256{,}00}
\newcommand{\PDosE}{372 965{,}00}
\newcommand{\PDosN}{8 572 251{,}00}
\newcommand{\PTresE}{372 911{,}00}
\newcommand{\PTresN}{8 573 249{,}00}
\newcommand{\PCuatroE}{373 011{,}00}
\newcommand{\PCuatroN}{8 573 255{,}00}

\newcommand{\FechaPago}{15 de agosto de 2018}
\newcommand{\MontoPago}{S/\ 20 000{,}00}
\newcommand{\MontoPagoLetras}{VEINTE MIL Y 00/100 SOLES}
\newcommand{\FormaPagoHistorica}{dinero en efectivo, sin medio de pago del sistema financiero}

% Nombres genéricos --- sustituir por los reales al ir a notaría
\newcommand{\NotarioNombre}{Dr.\ JUAN CARLOS RÍOS SALINAS (nombre genérico a sustituir)}
\newcommand{\NotariaCiudad}{Puerto Maldonado}
\newcommand{\ColegiaturaAbogado}{CAL Madre de Dios N.º 0000 (genérico)}
\newcommand{\AbogadoNombre}{Abog.\ ANA MARÍA QUISPE HUAMÁN (nombre genérico a sustituir)}
\newcommand{\IngenieroNombre}{Ing.\ PEDRO ANTONIO MAMANI CONDORI (nombre genérico a sustituir)}
\newcommand{\IngenieroCIP}{CIP 000000 (genérico)}
\newcommand{\Municipalidad}{Municipalidad Distrital de Inambari}
\newcommand{\Provincia}{Tambopata}
\newcommand{\Departamento}{Madre de Dios}
\newcommand{\Distrito}{Inambari}

\newcommand{\LugarOtorgamiento}{Puerto Maldonado}
\newcommand{\DiaOtorgamiento}{\faltante{[día]}}
\newcommand{\MesOtorgamiento}{\faltante{[mes]}}
\newcommand{\AnioOtorgamiento}{2026}

\newcommand{\TestigoUno}{\faltante{[testigo 1: nombres, DNI, domicilio]}}
\newcommand{\TestigoDos}{\faltante{[testigo 2: nombres, DNI, domicilio]}}

\InicioDocumento{Casuística del remanente $\sim$90 ha: compraventa vs donación a Lucio}{DOC-25}

\noindent Este cuadro es de \textbf{otro expediente} (el de Lucio). El de Nico y Lucy cubre solo las 10 ha. No se transfiere el remanente en la escritura de las 10 ha.

\vspace{0.2em}
\noindent \textbf{Hechos base.} Titular: \VendedorNombreCompleto, DNI \VendedorDNI, partida \PartidaMatriz, \AreaMatriz. Adquirió en 2017 por \PrecioAdquisicionMatriz. Ocupante del remanente: \IntermediarioNombre, DNI \IntermediarioDNI, sin título. Richar es soltero, sin unión de hecho vigente, y \textbf{tiene hijos}. El monto que Lucio habría entregado por las $\sim$90 ha \textbf{aún no está documentado}.

\Clausula{0. Condición previa (vale para los dos caminos)}
\textbf{Primero independizar las 10 ha} a favor de Nico y Lucy. Recién después Richar dispone del remanente. Si se ``vende o dona 90 ha'' sin recortar las 10, se arriesga la partida entera y el trato de 2018.

\section*{1. Cuadro comparativo --- compraventa vs donación}

{\small\renewcommand{\arraystretch}{1.22}
\noindent\begin{tabularx}{\textwidth}{>{\bfseries}p{3.15cm} >{\raggedright\arraybackslash}p{5.55cm} X}
\toprule
Tema & Compraventa (onerosa) & Donación (gratuita) \\
\midrule
Qué es & Richar vende el remanente a Lucio por un precio. & Richar regala el remanente a Lucio. \\
Cuándo encaja & Hubo o habrá dinero (aunque sea informal). Es lo que las partes han descrito: Lucio ``está comprando lo que queda''. & Solo si \textbf{nunca} hubo precio ni se pagará. No encaja si ya pasó plata. \\
Documento & Minuta + escritura de compraventa. Precio, forma de pago, posesión. & Escritura pública de donación con aceptación de Lucio (art.\ 1625 CC). Papel privado no vale. \\
Independización & Igual: plano UTM, recorte 10 ha, remanente $\sim$90 ha. & Igual. La donación no salta el plano ni SUNARP. \\
Predial atrasado & Lo liquida y paga \textbf{Richar} (titular al 1 de enero). Constancia de no adeudo. & Igual: Richar. \\
Alcabala & La paga \textbf{Lucio}. 3\,\% sobre el mayor entre precio y autoavalúo de las $\sim$90 ha, menos 10 UIT (S/\ 55\,000 en 2026). Las 90 ha tienen más chance de superar ese umbral que las 10 ha. & \textbf{También la paga Lucio.} La donación de inmueble \textbf{sí grava alcabala} (base: autoavalúo del remanente). No se ahorra Mazuko. \\
Medio de pago (Ley 28194) & Si el precio $>$ 3 UIT (S/\ 16\,500 en 2026): depósito o transferencia a cuenta de \textbf{Richar}. Conservar vouchers (como los S/\ 38\,000 de 2016 a Saturnino). & No hay precio que bancarizar. Por eso a veces se pretende donar: es el camino equivocado si sí hubo dinero. \\
Notaría y SUNARP & Escritura + inscripción del remanente a nombre de Lucio. Lo habitual: los paga Lucio. & Igual. No hay atajo registral. \\
Hijos de Richar & Una venta con precio real es más defendible. Los hijos \textbf{no firman} la compraventa onerosa. & Donar $\sim$90\,\% de este fundo es el escenario de \textbf{mayor riesgo}: pueden pedir reducción por afectar la legítima (porción disponible $\approx$ 1/3 si hay descendientes). \\
Los S/\ 20\,000 de las 10 ha & Ajenos. No se mezclan. & Ajenos. No se ``convierten'' en donación de 90 ha. \\
Posesión de Lucio & Ya ocupa el remanente. La escritura declara esa posesión. No basta para ser dueño. & Igual: ocupar no dona ni vende. \\
Cónyuge / conviviente de Lucio & Si Lucio es casado o conviviente, casi seguro interviene en la adquisición. Dato aún faltante. & Igual. \\
Simulación & Riesgo si el precio es irrisorio respecto del autoavalúo. & Riesgo máximo si se etiqueta ``donación'' un trato que fue compra. \\
Recomendación en este fundo & \textbf{Camino propio} si Lucio pagó o va a pagar. & \textbf{No recomendada} aquí, salvo liberalidad pura y dictamen sucesorio. \\
\bottomrule
\end{tabularx}}

\section*{2. Casuísticas (qué hacer según los hechos)}

{\small\renewcommand{\arraystretch}{1.2}
\noindent\begin{tabularx}{\textwidth}{>{\bfseries}c p{4.35cm} X}
\toprule
N.º & Hecho & Camino &
\midrule
1 & Lucio \textbf{ya pagó} a Richar (monto a documentar) por ``el resto''. & \textbf{Compraventa} del remanente. Reconocer el pago histórico en la escritura. \textbf{No} donar. No volver a pagar el mismo monto. Si en su día superó 3 UIT y fue efectivo, el notario deja constancia (como en las 10 ha); no se ``arregla'' depositando ahora la misma suma. \\
2 & Lucio \textbf{pagará ahora} más de S/\ 16\,500. & \textbf{Compraventa}. Depósito o transferencia a la cuenta de Richar \textbf{antes} de la escritura. Insertar vouchers. Liquidar alcabala. \\
3 & Lucio pagará ahora \textbf{S/\ 16\,500 o menos}. & Compraventa. El umbral de medio de pago no obliga, pero el depósito sigue siendo más limpio. \\
4 & Nadie recuerda el monto; hay trato de compra y Lucio ocupa. & Averiguar y \textbf{escribir el precio real}. Si no hay precio, no hay compraventa; entonces el único acto honesto sería donación, con el riesgo sucesorio de la casuística 6. \\
5 & Quieren donación \textbf{para evitar} alcabala, notario o SUNARP. & \textbf{No funciona.} Donación también exige escritura, alcabala, predial y registro. \\
6 & Richar dona $\sim$90 ha teniendo \textbf{hijos}. & Muy riesgosa. 90 ha de 100 es casi todo este bien. Los hijos pueden atacar. Mejor venta por precio. \\
7 & Donación de 90 ha y, aparte, Nico/Lucy compran 10 ha. & Dos actos distintos. No usar la donación de Lucio para ``cubrir'' las 10 ha, ni al revés. \\
8 & Vender a Lucio las \textbf{100 ha} y que él luego pase 10 ha a Nico/Lucy. & \textbf{Peor.} Doble alcabala, Lucio sería dueño formal de las 10 ha, y se borra el trato directo de 2018. No. \\
9 & Lucio casado o con unión de hecho. & Su cónyuge/conviviente interviene. Completar estado civil (hoy está en blanco). \\
10 & Independizar 10 ha y dejar el remanente \textbf{a nombre de Richar} años más. & Válido. Lucio sigue de poseedor informal hasta su propia escritura. El riesgo es de Lucio y de Richar, no un título de Nico/Lucy. \\
\bottomrule
\end{tabularx}}

\section*{3. Orden práctico recomendado}

\begin{enumerate}[leftmargin=1.15cm,itemsep=0.2em]
\item Cerrar \textbf{primero} las 10 ha (este expediente): plano, predial de Richar, alcabala Nico/Lucy, escritura, SUNARP.
\item Lucio firma el DOC-16 (no reclama las 10 ha). Eso \textbf{no} le da las 90 ha.
\item Abrir expediente de Lucio: DNI, estado civil, domicilio, \textbf{cuánto pagó o pagará}, cuenta de Richar.
\item Plano del remanente (sale del mismo levantamiento de independización).
\item Predial al día sobre lo que quede en la partida de Richar.
\item \textbf{Minuta de compraventa} Richar $\rightarrow$ Lucio del remanente (no donación, si hubo o habrá precio).
\item Alcabala de Lucio en Mazuko + escritura + inscripción.
\end{enumerate}

\vspace{0.25em}
\noindent Pro-rata solo de referencia (no es el precio de Lucio): Richar pagó \PrecioAdquisicionMatriz\ por 100 ha en 2017 $\rightarrow$ 90 ha $\approx$ S/\ 58\,500 a valores de entonces. El precio de Lucio puede ser otro; hay que declararlo. El municipio puede tomar un autoavalúo mayor para alcabala.

\LugarFecha
\noindent
\TresFirmas{EL TITULAR}{\VendedorNombreCompleto\\DNI \VendedorDNI}{LUCIO (REMANENTE)}{\IntermediarioNombre\\DNI \IntermediarioDNI}{ENTERADOS 10 ha}{\CompradorNombreCompleto}

\end{document}
